\documentclass[11pt,twoside,openany]{book}

% ------------------------------------------------
% Encodage et langue
% ------------------------------------------------

% Encodage : avec fontspec (XeLaTeX/LuaLaTeX), UTF-8 est natif ; ne pas charger inputenc/fontenc
\usepackage[french]{babel}

% ------------------------------------------------
% Polices (nécessite XeLaTeX ou LuaLaTeX)
% ------------------------------------------------

\usepackage{fontspec}
\setmainfont{Times New Roman}

% ------------------------------------------------
% Template couverture
% ------------------------------------------------

\usepackage{stys/Titlepage}

% ------------------------------------------------
% Couleurs
% ------------------------------------------------

\usepackage[dvipsnames,svgnames,x11names,table]{xcolor}
\definecolor{nuanbai}{HTML}{F5F5F5}
\pagecolor{nuanbai!10}

% ------------------------------------------------
% Maths
% ------------------------------------------------

\usepackage{amsmath,amssymb}

% ------------------------------------------------
% Mise en page
% ------------------------------------------------

\usepackage[left=2cm,right=2cm,top=.8cm,bottom=3.3cm]{geometry}
\setlength{\parindent}{0pt}

\usepackage{xpatch}
\usepackage[automark]{scrlayer-scrpage}
\usepackage[explicit,clearempty,pagestyles,newparttoc]{titlesec}
\usepackage{titletoc}
\titleformat{\chapter}[display]{\normalfont\fontsize{18}{22}\selectfont\bfseries}{}{0pt}{#1}
\titleformat{\section}{\normalfont\fontsize{14}{18}\selectfont\bfseries}{}{0pt}{#1}
\titleformat{\subsection}{\normalfont\fontsize{12}{16}\selectfont\bfseries}{}{0pt}{#1}
\titleformat{\subsubsection}{\normalfont\fontsize{11}{14}\selectfont\bfseries}{}{0pt}{#1}

\setcounter{secnumdepth}{5}
\setcounter{tocdepth}{2}

% ------------------------------------------------
% Packages graphiques
% ------------------------------------------------

\usepackage{hyperref}
\hypersetup{colorlinks=true,linkcolor=black,urlcolor=blue,pdfborder={0 0 0}}
\usepackage{fontawesome5}

\usepackage{tikz}
\usetikzlibrary{calc,shadows}

\usepackage[most]{tcolorbox}
\tcbuselibrary{breakable,skins,theorems}

\usepackage{tabularx}
\usepackage{lastpage}
\usepackage{varwidth}
\usepackage{colortbl}
\usepackage{tabularray}
\usepackage{pgfornament}
\usepackage{adjustbox}
\usepackage{graphicx}

\usepackage{pgfplots}
\pgfplotsset{compat=1.18}

% ------------------------------------------------
% Texte de démonstration
% ------------------------------------------------

\usepackage{lipsum}

% ------------------------------------------------
% Couleurs personnalisées
% ------------------------------------------------

\definecolor{qing}{HTML}{26A69A}

% ------------------------------------------------
% Environnements
% ------------------------------------------------

\usepackage{amsthm}
\newtheorem*{solution}{\textbf{Solution :}}

% ------------------------------------------------
% Commande exercice
% ------------------------------------------------

\newcommand{\exercise}[2][\bullet]{\bigskip
\begin{tikzpicture}[remember picture,overlay]
\draw[line width=2pt,loosely dotted,teal]
(0,0)--node[pos=0.4,
rectangle,
minimum height=1.5em,
font=\sffamily\Large,
text=black,
fill=black!2,
drop shadow={opacity=.3, shadow xshift=0.1cm},
inner sep=1.5mm,
anchor=west] {$#1$  ~ #2} (\linewidth,0);
\end{tikzpicture}\bigskip\smallskip
}

% ------------------------------------------------
% DOCUMENT
% ------------------------------------------------

\begin{document}

\title{Projet d'application graphique pour l'édition d'un dictionnaire multilingue}
\author{}
\subtitle{Documentation technique}
\edition{Version 1.0}
\bookseries{Universite Sorbonne Paris-Nord}
\pressname{Sup Galilée}
\presslogo{figures/sup-galilee-logo.jpg}

\makecover

\tableofcontents
\newpage

% ------------------------------------------------
% Contenu
% ------------------------------------------------

\chapter{Introduction}

\section{Présentation du projet}

L'Éditeur de dictionnaires multilingues est une application de bureau développée en JavaFX, entièrement dédiée à la création, à la modification et à la gestion de dictionnaires lexicographiques au format LIFT (Lexicon Interchange Format). Conçue pour répondre aux besoins des linguistes et des lexicographes, cette application permet de parcourir, modifier et enrichir des dictionnaires multilingues avec un support complet et flexible des langues objet (formes lexicales, exemples) et des méta-langues (définitions, traductions, glosses).

Le format LIFT constitue un standard XML reconnu dans le domaine de la lexicographie, notamment pour les projets de documentation des langues. L'éditeur s'inscrit dans cette démarche en offrant une interface graphique intuitive qui facilite la saisie et l'organisation des données lexicales, tout en garantissant la conformité avec la spécification officielle.

L'application repose sur une architecture modulaire en deux parties : le module \texttt{lift-api}, qui assure la lecture et l'écriture des fichiers LIFT ainsi que la manipulation du modèle de données, et le module \texttt{dictionary-editor-fx}, qui fournit l'interface graphique complète. L'ensemble est développé en Java 21 et s'appuie sur Maven 3.x pour la construction et la gestion des dépendances.

\section{Objectifs}

Les objectifs principaux de l'éditeur sont les suivants :

\subsection{Objectifs fonctionnels}

\begin{itemize}
\item Créer et modifier des dictionnaires multilingues au format LIFT standard, en respectant la structure XML définie par la spécification officielle.
\item Supporter les langues objet et méta-langues de manière flexible : l'utilisateur peut configurer les langues utilisées dans le dictionnaire et les ajouter ou les supprimer selon les besoins du projet.
\item Gérer l'ensemble des éléments lexicographiques : entrées, sens, exemples, variantes, prononciations, étymologies, relations, traits, annotations et champs personnalisés.
\item Exporter les données en CSV selon la vue courante affichée à l'écran, avec des en-têtes de colonnes correspondant exactement au tableau et un formatage conforme à la norme RFC~4180.
\item Valider les dictionnaires et assurer la cohérence des données, notamment via les outils de validation intégrés.
\item Offrir des fonctionnalités d'annulation et de rétablissement (undo/redo) pour faciliter le travail d'édition.
\end{itemize}

\subsection{Objectifs techniques}

\begin{itemize}
\item Maintenir une séparation nette entre la couche de données (\texttt{lift-api}) et l'interface (\texttt{dictionary-editor-fx}), afin de permettre la réutilisation du module API dans d'autres contextes (scripts, services, applications web).
\item Assurer la conformité avec la spécification LIFT officielle et la compatibilité avec les outils du domaine (FieldWorks, WeSay, FLEx).
\item Fournir une interface réactive et ergonomique basée sur JavaFX, avec support du multilingue et de l'internationalisation (français, anglais).
\end{itemize}

\section{Public cible}

Cette documentation s'adresse à plusieurs types de lecteurs. Les linguistes et lexicographes y trouveront les informations nécessaires pour utiliser l'éditeur dans le cadre de projets de documentation linguistique ou de création de dictionnaires. Les développeurs souhaitant contribuer au projet ou l'adapter à des besoins spécifiques disposeront d'une description détaillée de l'architecture et des composants techniques. Enfin, les utilisateurs occasionnels pourront s'appuyer sur les chapitres consacrés à l'installation et aux fonctionnalités pour prendre en main l'application rapidement.

\section{Contexte et cas d'usage}

L'éditeur trouve son utilité dans plusieurs contextes. Les projets de documentation des langues (language documentation) nécessitent souvent la création de dictionnaires bilingues ou multilingues pour accompagner la description grammaticale et les corpus de textes. Les lexicographes travaillant sur des langues peu dotées ou en danger peuvent utiliser l'outil pour structurer et enrichir progressivement leurs données. Les enseignants et chercheurs en linguistique y trouvent un moyen de constituer des lexiques spécialisés (terminologie, dialectologie). Enfin, l'éditeur peut servir de base pour des projets d'édition collaborative, dans la mesure où le format LIFT est interopérable avec d'autres outils du domaine (FieldWorks, WeSay, etc.).

\section{Documentation en ligne}

Une documentation technique complémentaire, régulièrement mise à jour, est disponible en ligne à l'adresse suivante :

\url{https://multilingual-dictionary-editor.vercel.app/}

Cette documentation web propose une navigation structurée par thèmes (Introduction, Architecture, Installation, Format LIFT, Fonctionnalités, Glossaire) ainsi qu'une barre de recherche intégrée permettant de retrouver rapidement un terme, une page ou un concept. Elle constitue un complément utile au présent document pour les utilisateurs souhaitant approfondir certains aspects techniques.
\section{Plan du document}

Le présent document est organisé en six chapitres. Après cette introduction, le chapitre~2 décrit l'architecture et la structure du projet (modules Maven, organisation du code, technologies, flux de données). Le chapitre~3 présente le format LIFT en détail (structure, éléments clés, ranges, fichiers associés, validation). Le chapitre~4 fournit les instructions d'installation et d'utilisation (prérequis, compilation, lancement, tests, dépannage). Le chapitre~5 détaille les fonctionnalités de l'interface (disposition, navigation, tableaux, formulaire d'édition, configuration, outils, internationalisation). Enfin, le chapitre~6 présente un glossaire des termes techniques utilisés dans l'éditeur et la documentation LIFT.


\chapter{Architecture et structure du projet}

\section{Organisation Maven}

Le projet adopte une structure Maven multi-modules, ce qui permet de séparer clairement les responsabilités entre la couche de données et l'interface utilisateur. Le fichier \texttt{pom.xml} situé à la racine du dépôt agit comme agrégateur et déclare les deux modules principaux. Cette organisation facilite la maintenance du code, les tests unitaires ciblés et la réutilisation éventuelle du module \texttt{lift-api} dans d'autres projets.

L'arborescence du dépôt est la suivante :

\begin{verbatim}
multilingual-dictionary-editor/
├── pom.xml              (agrégateur)
├── lift-api/            (module bibliothèque)
└── dictionary-editor-fx/ (module application)
\end{verbatim}

\subsection{Schéma d'architecture}

Le schéma ci-dessous illustre les relations entre les différents composants. Le module \texttt{dictionary-editor-fx} dépend du module \texttt{lift-api} pour accéder au modèle de données et aux fonctionnalités de chargement et de sauvegarde des fichiers LIFT.

\begin{center}
\begin{tikzpicture}[
  box/.style={rectangle, draw=qing!80, fill=qing!10, rounded corners=3pt, minimum width=3.2cm, minimum height=0.7cm},
  arrow/.style={->, >=stealth, thick}
]
  \node[box] (root) at (0,0) {\texttt{pom.xml (agrégateur)}};
  \node[box] (api) at (-3.2,-1.5) {\texttt{lift-api}};
  \node[box] (fx) at (3.2,-1.5) {\texttt{dictionary-editor-fx}};
  \draw[arrow] (root) -- (api);
  \draw[arrow] (root) -- (fx);
  \draw[arrow] (fx.south) to[bend right=25] node[below, font=\scriptsize] {dépend de} (api.south);
\end{tikzpicture}
\end{center}

\section{Module lift-api}

Le module \texttt{lift-api} constitue le cœur technique du projet. Il contient l'ensemble du modèle de données correspondant à la structure LIFT, ainsi que les parsers et serializers XML permettant de charger et de sauvegarder les dictionnaires. Ce module est conçu pour être réutilisable : il peut être intégré dans d'autres applications Java sans dépendre de JavaFX.

Les classes principales fournies par ce module sont les suivantes :

\begin{itemize}
\item \texttt{LiftDictionary} : représente le dictionnaire complet ; gère le chargement depuis un fichier, la sauvegarde et l'accès aux composants (entrées, header, etc.).
\item \texttt{LiftDictionaryLoader} : responsable du chargement et du parsing des fichiers LIFT.
\item \texttt{LiftEntry}, \texttt{LiftSense}, \texttt{LiftExample} : entités lexicales de base, organisées hiérarchiquement.
\item \texttt{MultiText}, \texttt{Form} : gestion du contenu multilingue (formes, définitions, glosses, etc.) avec support des propriétés JavaFX pour le binding.
\item \texttt{LiftTrait}, \texttt{LiftAnnotation}, \texttt{LiftField} : métadonnées et champs personnalisés attachés aux entrées, sens ou exemples.
\item \texttt{LiftVariant}, \texttt{LiftPronunciation}, \texttt{LiftEtymology}, \texttt{LiftRelation} : éléments associés aux entrées (variantes orthographiques, prononciations, étymologies, relations sémantiques).
\end{itemize}

Les tests unitaires, situés dans \texttt{src/test/java/}, vérifient le bon fonctionnement du chargement et de la manipulation des données. Des fichiers LIFT d'exemple sont fournis dans \texttt{src/test/resources/lift/} pour alimenter ces tests.

\section{Module dictionary-editor-fx}

Le module \texttt{dictionary-editor-fx} regroupe l'ensemble de l'interface graphique. Il s'appuie sur JavaFX pour offrir une application de bureau réactive et ergonomique. L'architecture suit un modèle de type MVC (Model-View-Controller) : le \texttt{MainController} orchestre les vues et les interactions, tandis que le modèle est fourni par \texttt{lift-api}.

Les composants principaux sont les suivants :

\begin{itemize}
\item \texttt{MainApp} : point d'entrée de l'application ; initialise la fenêtre principale et charge la vue FXML.
\item \texttt{MainController} : contrôleur principal ; gère la navigation entre les vues (Entrées, Sens, Exemples, etc.), le chargement des tableaux dynamiques, le formulaire d'édition et les actions utilisateur (ouverture, sauvegarde, export CSV, etc.).
\item \texttt{MainView.fxml} : fichier FXML définissant la mise en page de l'interface (arbre de navigation à gauche, zone centrale pour les tableaux, panneau d'édition à droite).
\item \texttt{controls/} : package contenant les éditeurs spécialisés (\texttt{SenseEditor}, \texttt{ExampleEditor}, \texttt{MultiTextEditor}, \texttt{NotableEditor}, \texttt{VariantEditor}, \texttt{RelationEditor}, \texttt{EntryEditor}, etc.), chacun étant responsable de l'affichage et de l'édition d'un type d'élément précis. Le \texttt{MultiTextEditor} gère les champs multilingue avec un onglet ou un champ par langue ; les éditeurs de notes, variantes et relations réutilisent des composants communs pour assurer une cohérence visuelle.
\item \texttt{undo/} : implémentation du pattern Command pour l'annulation et le rétablissement des modifications (\texttt{DeleteEntryCommand}, \texttt{DeleteSenseCommand}, \texttt{DeleteExampleCommand}, etc.).
\item \texttt{core/} : services et exceptions (\texttt{DictionaryService}, \texttt{LiftOpenException}) pour la gestion des fichiers et des erreurs.
\end{itemize}

\section{Arborescence détaillée des sources}

Le module \texttt{lift-api} organise son code source dans \texttt{src/main/java/fr/cnrs/lacito/liftapi/}, avec les sous-packages \texttt{model/} (entités), \texttt{xml/} (parsers, \texttt{LiftSaxHandler}, \texttt{LiftWriter}) et \texttt{exporter/} (export LaTeX notamment). Le module \texttt{dictionary-editor-fx} structure son code dans \texttt{src/main/java/fr/cnrs/lacito/liftgui/} avec \texttt{ui/} (contrôleur, vues FXML, contrôles personnalisés), \texttt{core/}, \texttt{data/} et \texttt{undo/}. Les ressources (fichiers FXML, feuilles de style CSS, fichiers de traduction \texttt{messages\_fr.properties} et \texttt{messages\_en.properties}) se trouvent dans \texttt{src/main/resources/}.

\section{Technologies utilisées}

Le projet utilise un stack technique moderne et stable :

\begin{itemize}
\item \textbf{Java 21} : le code est compilé avec l'option \texttt{release 21}, ce qui garantit la compatibilité avec les fonctionnalités récentes du langage (records, pattern matching, virtual threads) tout en assurant une exécution optimale. Le modèle de données utilise les propriétés JavaFX (\texttt{StringProperty}, \texttt{ListProperty}) pour le binding bidirectionnel avec l'interface.
\item \textbf{JavaFX 21} : les modules \texttt{javafx-controls} et \texttt{javafx-fxml} fournissent les composants d'interface (TableView, TextField, TitledPane, TreeView, etc.) et le système de mise en page déclarative. L'interface est entièrement construite en JavaFX, sans dépendance à Swing.
\item \textbf{Maven 3.x} : outil de construction et de gestion des dépendances ; permet une compilation reproductible, l'exécution des tests, le packaging et une intégration aisée dans les environnements de développement (Eclipse, IntelliJ IDEA) et de déploiement.
\item \textbf{Lombok} (optionnel, dans lift-api) : réduit le code boilerplate pour les getters, setters et constructeurs.
\end{itemize}

\section{Flux de données et performance}

Le chargement d'un dictionnaire LIFT s'effectue via un parser SAX (\texttt{LiftSaxHandler}) qui construit progressivement le modèle en mémoire. Pour les dictionnaires volumineux, le parsing reste efficace car il ne charge pas l'intégralité du document en mémoire. Lors de l'affichage des vues Sens ou Exemples, une table de correspondance (map) précalculée associe chaque sens ou exemple à son entrée parente, évitant des recherches linéaires répétées et garantissant des performances O(1) pour la navigation. L'export CSV s'appuie sur la vue courante (\texttt{getCurrentTableView()}) pour déterminer le tableau à exporter, ce qui assure la cohérence entre l'affichage et le fichier généré.


\chapter{Format LIFT}

\section{Présentation}

LIFT (Lexicon Interchange Format) est un format XML standard conçu spécifiquement pour les dictionnaires lexicographiques. Développé à l'origine dans le cadre des projets de documentation des langues (notamment par SIL International), il est aujourd'hui largement utilisé dans le domaine de la lexicographie et de la linguistique de terrain. Le format est maintenu par la communauté SIL et documenté sur le dépôt \texttt{sillsdev/lift-standard}.

Le format LIFT offre un support complet du multilingue : les entrées lexicales peuvent comporter des formes dans plusieurs langues objet, tandis que les définitions, glosses et traductions utilisent les méta-langues. Il intègre également de nombreux éléments d'enrichissement : traits (paires nom/valeur), annotations, notes, relations sémantiques entre entrées, champs personnalisés, prononciations, étymologies, etc. Cette richesse structurelle permet de modéliser des dictionnaires complexes tout en restant conforme à une spécification bien définie. Les versions 0.13 et 0.15 de la spécification sont supportées par l'éditeur.

\section{Structure principale}

Un fichier LIFT est un document XML dont la structure de base comprend deux parties principales :

\begin{itemize}
\item \texttt{header} : contient les métadonnées du dictionnaire, la définition des ranges (taxinomies hiérarchiques, par exemple pour les domaines sémantiques), les field-definitions (définitions des champs personnalisés), ainsi que la déclaration des langues objet et méta-langues utilisées.
\item \texttt{body} : contient la liste des éléments \texttt{entry}, chaque entrée représentant une unité lexicale (mot ou expression) avec l'ensemble de ses sens, exemples, variantes et éléments associés.
\end{itemize}

\subsection{Hiérarchie des éléments}

La hiérarchie des éléments au sein d'une entrée peut être schématisée ainsi. Une entrée (\texttt{entry}) contient typiquement une unité lexicale (\texttt{lexical-unit}), un ou plusieurs sens (\texttt{sense}), des variantes (\texttt{variant}), des prononciations (\texttt{pronunciation}) et des relations (\texttt{relation}). Chaque sens peut à son tour contenir une définition, une gloss, des exemples et des sous-sens.

\begin{center}
\begin{tikzpicture}[
  level 1/.style={sibling distance=3.5cm, level distance=1cm},
  edge from parent/.style={draw, ->},
  every node/.style={font=\small\ttfamily}
]
  \node {entry}
    child { node {lexical-unit} }
    child { node {sense}
      child { node {definition} }
      child { node {gloss} }
      child { node {example} }
    }
    child { node {variant} }
    child { node {pronunciation} }
    child { node {relation} };
\end{tikzpicture}
\end{center}

\section{Éléments clés}

Parmi les éléments les plus utilisés dans un dictionnaire LIFT, on retiendra :

\begin{itemize}
\item \texttt{lexical-unit} et \texttt{form} : la forme lexicale de l'entrée, exprimée dans une ou plusieurs langues objet. Chaque forme est associée à un code de langue.
\item \texttt{sense} : représente un sens de l'entrée. Il contient une \texttt{definition} et une \texttt{gloss} (souvent utilisée pour la traduction courte), rédigées dans les méta-langues.
\item \texttt{example} : phrase ou expression illustrant l'usage du sens. Les exemples peuvent comporter des traductions par type (par exemple «~traduction littérale~», «~traduction libre~»).
\item \texttt{trait} : paire nom/valeur permettant de catégoriser l'élément. Les traits les plus courants incluent \texttt{grammatical-info} (catégorie grammaticale : nom, verbe, adverbe, etc.) et \texttt{variant-type} pour les variantes.
\item \texttt{grammatical-info} : élément spécialisé pour la catégorie grammaticale, souvent lié à une taxinomie définie dans le header.
\item \texttt{note}, \texttt{relation}, \texttt{field} : respectivement les notes (avec type optionnel), les relations sémantiques vers d'autres entrées, et les champs personnalisés définis dans le header.
\end{itemize}

\section{Ranges et taxinomies}

Les ranges permettent de définir des taxinomies hiérarchiques réutilisables dans le dictionnaire. Un range typique est celui des domaines sémantiques (par exemple DDP, Dewey Decimal for Linguistics) : chaque élément du range possède un identifiant, un libellé multilingue, une description optionnelle et peut avoir un parent pour former une arborescence. Les traits et \texttt{grammatical-info} peuvent référencer ces ranges via l'attribut \texttt{href} pointant vers un fichier \texttt{.lift-ranges} externe. L'éditeur permet de consulter et modifier les ranges depuis la section Configuration.

\section{Fichiers associés}

Un projet LIFT complet comprend souvent plusieurs fichiers : le fichier principal \texttt{.lift} (contenant le header et le body), un ou plusieurs fichiers \texttt{.lift-ranges} pour les taxinomies, et éventuellement des fichiers de configuration. Les chemins relatifs sont utilisés pour référencer les fichiers ranges depuis le header. L'éditeur gère le chargement de ces fichiers associés lors de l'ouverture d'un dictionnaire.

\section{Exemple de structure XML minimale}

Un fichier LIFT minimal comporte une structure du type :

\begin{verbatim}
<?xml version="1.0" encoding="UTF-8"?>
<lift producer="..." version="0.13">
  <header>
    <ranges>...</ranges>
    <fields>...</fields>
  </header>
  <body>
    <entry id="...">
      <lexical-unit>
        <form lang="fr"><text>mot</text></form>
      </lexical-unit>
      <sense>
        <gloss lang="en"><text>word</text></gloss>
      </sense>
    </entry>
  </body>
</lift>
\end{verbatim}

Les versions supportées par l'éditeur incluent 0.13 et 0.15 de la spécification LIFT.

\section{Validation et conformité}

L'éditeur ne valide pas strictement le XML contre un schéma XSD au chargement ; il s'appuie sur un parser SAX tolérant qui construit le modèle à partir des éléments reconnus. Les éléments non reconnus ou mal formés peuvent être ignorés ou provoquer une erreur selon leur nature. L'outil de validation intégré (menu Outils) permet de vérifier la cohérence des données : comptage des entrées, sens, exemples, vérification des langues déclarées, détection d'éventuelles incohérences. Pour une validation XML stricte, il est recommandé d'utiliser des outils externes (xmllint, Oxygen XML) avec le schéma LIFT officiel.

\section{Ressources}

La spécification officielle du format LIFT, maintenue par la communauté SIL et les contributeurs du projet, est disponible sur le dépôt GitHub : \url{https://github.com/sillsdev/lift-standard}. Cette ressource constitue la référence pour toute implémentation ou extension du format. Des exemples de fichiers LIFT sont fournis dans le dépôt du projet, notamment dans \texttt{lift-api/src/test/resources/lift/} et \texttt{dictionary-editor-fx/src/main/resources/lift/}.


\chapter{Installation et utilisation}

\section{Prérequis}

Avant d'installer et de lancer l'éditeur, il est nécessaire de disposer des outils suivants sur la machine :

\begin{itemize}
\item \textbf{JDK 21} : le kit de développement Java version 21 doit être installé et configuré. Il est recommandé de définir la variable d'environnement \texttt{JAVA\_HOME} pour pointer vers le répertoire d'installation du JDK. Cette variable est utilisée par Maven et par de nombreux environnements de développement. Les distributions recommandées incluent OpenJDK, Eclipse Temurin, Oracle JDK ou Azul Zulu.
\item \textbf{Maven} : l'outil de construction Maven (version 3.6 ou supérieure) doit être installé et accessible depuis le chemin d'exécution du système. Maven permet de compiler le projet, de gérer les dépendances et de lancer l'application. Sur Windows, s'assurer que \texttt{MAVEN\_HOME} est configuré et que \texttt{mvn} est dans le \texttt{PATH}.
\item \textbf{JavaFX} : les dépendances JavaFX sont gérées automatiquement par Maven ; aucune installation manuelle n'est requise.
\end{itemize}

\subsection{Vérification de l'environnement}

Avant de lancer la compilation, on peut vérifier l'environnement avec les commandes suivantes :

\begin{verbatim}
java -version
javac -version
mvn -version
\end{verbatim}

La sortie doit indiquer une version Java 21 et Maven 3.6 ou supérieur.

\section{Compilation}

Pour obtenir une copie du projet et le compiler, exécuter les commandes suivantes depuis un terminal, en se plaçant dans le répertoire de travail souhaité :

\begin{verbatim}
git clone "https://github.com/I2S9/multilingual-dictionary-editor.git"
cd multilingual-dictionary-editor
mvn install
\end{verbatim}

La commande \texttt{mvn install} compile l'ensemble des modules et les installe dans le dépôt Maven local (généralement \texttt{\textasciitilde/.m2/repository/}). Cette étape est indispensable : le module \texttt{dictionary-editor-fx} dépend du module \texttt{lift-api}, et sans installation préalable de ce dernier, la compilation échouera avec une erreur de dépendance manquante.

\section{Lancement}

Une fois la compilation terminée, l'application peut être lancée avec la commande suivante :

\begin{verbatim}
mvn javafx:run -pl dictionary-editor-fx
\end{verbatim}

L'option \texttt{-pl dictionary-editor-fx} indique à Maven de cibler uniquement le module de l'interface graphique. Le plugin \texttt{javafx-maven-plugin} configure le point d'entrée et lance la fenêtre principale de l'éditeur.

En cas d'échec de la commande ci-dessus (par exemple en raison d'un problème de configuration de Maven ou du JDK), une alternative consiste à compiler et installer d'abord le module \texttt{lift-api}, puis à lancer l'éditeur en spécifiant explicitement le fichier \texttt{pom.xml} du module :

\begin{verbatim}
mvn install -pl lift-api
mvn -f dictionary-editor-fx/pom.xml javafx:run
\end{verbatim}

Sous Windows, un script batch \texttt{run-app.bat} est fourni à la racine du projet. Ce script configure l'environnement et lance l'application ; il peut être exécuté par double-clic ou depuis l'invite de commandes.

\subsection{Packaging et exécutable autonome}

Pour générer un JAR exécutable avec toutes les dépendances (fat JAR), on peut utiliser le plugin \texttt{maven-assembly-plugin} ou \texttt{maven-shade-plugin}. La commande \texttt{mvn package} produit un JAR dans \texttt{dictionary-editor-fx/target/}. Pour une distribution autonome, il est possible de configurer \texttt{jlink} ou \texttt{jpackage} afin de créer une image native ou un installateur pour Windows, macOS ou Linux.

\section{Tests}

Le projet inclut des tests unitaires pour le module \texttt{lift-api}. Pour exécuter l'ensemble des tests de tous les modules :

\begin{verbatim}
mvn test
\end{verbatim}

Pour exécuter uniquement les tests du module \texttt{lift-api} :

\begin{verbatim}
mvn test -pl lift-api
\end{verbatim}

Les tests vérifient notamment le chargement correct des fichiers LIFT, la manipulation du modèle de données et la cohérence des opérations de lecture et d'écriture. Les tests du module \texttt{lift-api} couvrent le parsing SAX, la sérialisation XML, les conversions de modèles et l'export LaTeX.

\section{Dépannage}

En cas de problème lors du lancement, vérifier les points suivants :

\begin{itemize}
\item \textbf{Erreur «~module not found~»} : s'assurer que le module \texttt{lift-api} a bien été installé avec \texttt{mvn install -pl lift-api} avant de lancer l'éditeur.
\item \textbf{Erreur JavaFX} : sur certaines configurations, JavaFX peut nécessiter des arguments de module (\texttt{--add-modules javafx.controls,javafx.fxml}). Le plugin \texttt{javafx-maven-plugin} les ajoute automatiquement.
\item \textbf{Fichier LIFT corrompu} : si un fichier LIFT ne s'ouvre pas, vérifier la validité XML (balises fermées, encodage UTF-8) et la conformité à la spécification LIFT.
\end{itemize}

\section{Notes Windows}

Sous Windows, il peut arriver que Maven affiche le message d'erreur «~La syntaxe de la commande n'est pas correcte~» lors de l'exécution de \texttt{mvn.cmd}. Cette erreur est souvent due à une configuration incorrecte de la variable \texttt{JAVA\_HOME} : si celle-ci contient un saut de ligne, un espace problématique ou un caractère parasite, le script Maven peut échouer. Il convient de vérifier que \texttt{JAVA\_HOME} pointe vers un chemin valide, sans caractères superflus. Sous PowerShell, on peut par exemple définir temporairement la variable avant de lancer Maven :

\begin{verbatim}
$env:JAVA_HOME = 'C:\Program Files\Java\jdk-21'
mvn -pl dictionary-editor-fx javafx:run
\end{verbatim}


\chapter{Fonctionnalités de l'interface}

\section{Disposition générale}

L'interface de l'éditeur est organisée selon un schéma en trois panneaux, offrant une séparation claire entre la navigation, la consultation des données et l'édition. Cette disposition permet à l'utilisateur de parcourir rapidement le contenu du dictionnaire tout en ayant un accès direct au formulaire de modification.

\begin{center}
\begin{tikzpicture}[
  pane/.style={rectangle, draw=qing!60, fill=qing!5, rounded corners=2pt, minimum height=1.2cm},
  label/.style={font=\small\ttfamily}
]
  \node[pane, minimum width=2.2cm] (nav) at (0,0) {\parbox{2cm}{\centering Navigation\\arbre}};
  \node[pane, minimum width=4.5cm] (table) at (4,0) {\parbox{4.2cm}{\centering Tableaux\\(Entrées, Sens, etc.)}};
  \node[pane, minimum width=2.8cm] (edit) at (8.5,0) {\parbox{2.5cm}{\centering Formulaire\\d'édition}};
  \draw[->, thin] (nav.east) -- (table.west);
  \draw[->, thin] (table.east) -- (edit.west);
\end{tikzpicture}
\end{center}

Le panneau de gauche contient l'arbre de navigation ; la zone centrale affiche les tableaux dynamiques dont le contenu varie selon la vue sélectionnée ; le panneau de droite présente le formulaire d'édition détaillé de l'élément sélectionné. Lorsqu'aucun élément n'est sélectionné, le panneau de droite affiche le message «~(sélectionner un élément)~».

\subsection{Barre de menu et raccourcis}

La barre de menu propose les actions principales : Fichier (Ouvrir, Enregistrer, Fermer), Édition (Annuler, Rétablir), Outils (Validation, Export CSV) et Aide. Les raccourcis clavier standards sont supportés : Ctrl+O pour ouvrir, Ctrl+S pour enregistrer, Ctrl+Z pour annuler, Ctrl+Y pour rétablir. L'export CSV est accessible depuis le menu Outils ou via un bouton dédié dans la barre d'outils.

\section{Navigation}

L'arbre de navigation, situé dans le panneau gauche, structure l'accès aux différentes vues du dictionnaire. Il est organisé en plusieurs groupes :

\begin{itemize}
\item \textbf{Objets} : regroupe les vues par type d'élément lexical : Entrées, Sens, Exemples, Notes, Variantes, Étymologies et Relations. Chaque vue affiche un tableau listant les éléments correspondants avec des colonnes adaptées (formes, glosses, types, etc.).
\item \textbf{Langues} : propose deux vues dédiées aux champs multilingues : Langues objet (formes lexicales, exemples, prononciations) et Méta-langues (définitions, glosses, traductions). Ces vues permettent de parcourir l'ensemble des contenus textuels par langue.
\item \textbf{Category browser} : offre des vues par catégorie : Grammatical info, Traits, Annotations, Types de notes, Types de relations, Types de traductions, Fields. Chaque vue affiche les valeurs distinctes avec leur fréquence d'occurrence.
\item \textbf{Configuration} : accès aux paramètres du header du dictionnaire (Description, Gestion des langues, Ranges, définitions de champs, etc.).
\item \textbf{Saisie rapide} : mode de création d'entrées en lot, permettant de saisir plusieurs formes lexicales et informations grammaticales avant de les créer en une seule opération.
\end{itemize}

\subsection{Tableaux et filtres}

Chaque vue de tableau (Entrées, Sens, Exemples, etc.) affiche des colonnes adaptées au type d'objet. Les colonnes de texte multilingue (formes, glosses, définitions) disposent d'un filtre texte permettant de rechercher rapidement dans les données. Les filtres sont appliqués en temps réel ; seules les lignes correspondant aux critères sont affichées. L'export CSV prend en compte ces filtres : seules les lignes visibles sont exportées. Les colonnes peuvent être réordonnées et redimensionnées selon les préférences de l'utilisateur.

\subsection{Lien vers l'objet parent}

Dans les vues Sens, Exemples, Variantes, etc., un champ «~parent~» ou «~entrée~» affiche un lien cliquable vers l'entrée ou le sens parent. Un clic sur ce lien navigue directement vers l'objet parent et sélectionne l'élément correspondant, facilitant la navigation dans la hiérarchie du dictionnaire.

\section{Category browser}

Le Category browser constitue un outil puissant pour explorer le dictionnaire par catégories. Plutôt que de parcourir les entrées une à une, l'utilisateur peut consulter les valeurs distinctes d'une catégorie donnée (par exemple toutes les catégories grammaticales utilisées, ou tous les noms de traits) avec le nombre d'occurrences de chaque valeur. Un bouton «~Accéder aux objets correspondants~» permet d'afficher dans un tableau temporaire les entrées, sens ou exemples qui possèdent la valeur sélectionnée, facilitant ainsi la navigation ciblée et la vérification des données.

Les vues disponibles incluent : Grammatical info (catégories grammaticales), Traits (paires nom/valeur avec fusion des lignes identiques), Annotations, Types de notes, Types de relations, Types de traductions, et Fields (champs personnalisés). Chaque valeur affichée est accompagnée du nombre d'occurrences dans le dictionnaire, ce qui permet d'identifier rapidement les catégories les plus utilisées ou les valeurs orphelines.

\section{Formulaire d'édition}

Le panneau d'édition à droite affiche un formulaire contextuel selon le type d'élément sélectionné. Pour une entrée, on trouve les champs : unité lexicale (formes par langue), variantes, prononciations, étymologies, relations, notes et champs personnalisés. Pour un sens : définition, gloss, exemples, traits (dont grammatical-info), notes et champs. Pour un exemple : texte source, traductions par type, notes. Les champs multilingue utilisent le composant \texttt{MultiTextEditor}, qui affiche un onglet ou un champ par langue configurée. Les titres de section du formulaire (par exemple «~Définition~», «~Gloss~», «~Exemples~») sont stylés avec la classe \texttt{.editor-section-title} pour une hiérarchie visuelle claire.

\section{Configuration du dictionnaire}

La section Configuration permet de gérer les paramètres globaux du dictionnaire. Les langues objet et méta-langues peuvent être ajoutées ou supprimées ; les langues objet sont utilisées pour les formes lexicales (lexical-unit, exemples, prononciations), tandis que les méta-langues servent aux définitions, glosses et traductions. La configuration inclut également la gestion des types de notes, des types de traductions, des types d'annotations et des types de relations. Les définitions de champs et de traits (New definition) permettent de créer des champs personnalisés réutilisables dans les entrées. Enfin, les ranges (taxinomies hiérarchiques) définis dans le header peuvent être consultés et modifiés depuis cette section.

\section{Outils}

L'éditeur intègre plusieurs outils complémentaires :

\begin{itemize}
\item \textbf{Validation} : permet de vérifier la cohérence du dictionnaire et d'afficher un résumé des éléments (nombre d'entrées, de sens, d'exemples, langues utilisées). Utile pour détecter d'éventuelles incohérences avant export ou partage.
\item \textbf{Export CSV} : exporte les données de la vue courante au format CSV. Selon la page affichée (Entrées, Sens, Exemples, etc.), le tableau correspondant est exporté avec des en-têtes de colonnes identiques à ceux de l'interface. Le format respecte la norme RFC~4180 (séparateur virgule, échappement des guillemets et des retours à la ligne). Les filtres appliqués à l'écran sont pris en compte : seules les lignes visibles sont exportées.
\item \textbf{Undo / Redo} : annulation et rétablissement des modifications. Les suppressions d'entrées, de sens ou d'exemples peuvent être annulées ; les raccourcis clavier Ctrl+Z et Ctrl+Y sont disponibles. Le système utilise le pattern Command : chaque action destructrice (\texttt{DeleteEntryCommand}, \texttt{DeleteSenseCommand}, \texttt{DeleteExampleCommand}) est encapsulée dans une commande réversible.
\end{itemize}

\section{Internationalisation}

L'éditeur supporte l'internationalisation (i18n) via des fichiers de propriétés : \texttt{messages\_fr.properties} pour le français et \texttt{messages\_en.properties} pour l'anglais. Les libellés de l'interface (menus, boutons, messages d'erreur, titres de colonnes) sont externalisés et chargés selon la locale du système. Les clés de traduction suivent la convention \texttt{editor.*} pour les éléments de l'éditeur et \texttt{section.*} pour les titres de navigation.


\chapter{Glossaire}

Le glossaire ci-dessous regroupe les principaux termes techniques utilisés dans l'éditeur et dans la documentation du format LIFT. Ces définitions permettent de clarifier le vocabulaire spécifique au domaine de la lexicographie et au format de données.

\begin{description}
\item[LIFT] Lexicon Interchange Format. Format XML standard pour les dictionnaires lexicographiques, développé notamment dans le cadre des projets de documentation des langues. Il définit une structure hiérarchique pour les entrées, sens, exemples et métadonnées associées.

\item[Entrée (Entry)] Unité lexicale de base du dictionnaire. Une entrée représente un mot ou une expression et contient l'ensemble de ses sens, variantes orthographiques, prononciations, étymologies, relations et métadonnées. Chaque entrée possède un identifiant unique.

\item[Sens (Sense)] Définition ou traduction d'une entrée dans une ou plusieurs méta-langues. Un sens peut contenir une définition textuelle, une gloss (traduction courte), des exemples, des sous-sens, des relations sémantiques et des traits (par exemple la catégorie grammaticale).

\item[Exemple (Example)] Phrase ou expression illustrant l'usage d'un sens. Un exemple comporte généralement un texte dans la langue objet et des traductions par type (traduction littérale, traduction libre, etc.) dans les méta-langues.

\item[Trait (Trait)] Paire nom/valeur permettant de catégoriser un élément (entrée, sens, exemple, variante). Les traits les plus courants incluent \texttt{grammatical-info} (catégorie grammaticale) et \texttt{variant-type} pour les variantes orthographiques.

\item[Annotation] Métadonnée attachée à un élément du dictionnaire. Une annotation peut contenir des informations sur l'auteur (qui), la date (quand), ou d'autres attributs définis dans le header.

\item[Langue objet] Langue du dictionnaire, c'est-à-dire la langue des entrées lexicales. Elle est utilisée pour les formes (lexical-unit), les exemples et les prononciations. Un dictionnaire peut comporter plusieurs langues objet.

\item[Méta-langue] Langue utilisée pour les définitions, glosses et traductions. Les méta-langues permettent de décrire le sens des entrées dans une ou plusieurs langues de description. Un dictionnaire peut comporter plusieurs méta-langues.

\item[Category browser] Vue par catégories permettant de parcourir les valeurs distinctes d'une catégorie donnée (traits, types de notes, etc.) avec leur fréquence d'occurrence. Le Category browser offre un accès direct aux objets (entrées, sens, exemples) correspondant à une valeur sélectionnée.

\item[Range] Taxinomie hiérarchique définie dans le header du dictionnaire ou dans un fichier \texttt{.lift-ranges} externe. Un range permet de structurer des valeurs contrôlées (domaines sémantiques, catégories grammaticales) avec des relations parent-enfant. Les éléments du range possèdent un identifiant, un libellé multilingue et une description optionnelle.

\item[Field (champ personnalisé)] Champ défini dans le header via \texttt{field-definition} et pouvant être attaché aux entrées, sens ou exemples. Les champs personnalisés permettent d'ajouter des métadonnées spécifiques au projet sans modifier la structure LIFT standard.

\item[RFC 4180] Norme définissant le format CSV (Comma-Separated Values). L'éditeur respecte cette norme pour l'export : séparateur virgule, échappement des guillemets par double guillemet, encadrement des champs contenant virgules ou retours à la ligne. Un BOM UTF-8 est ajouté en tête du fichier pour une reconnaissance correcte de l'encodage.

\item[FXML] Format XML de mise en page pour JavaFX. Les fichiers \texttt{.fxml} définissent la structure des vues (hiérarchie des composants, liaisons avec le contrôleur) de manière déclarative, séparant ainsi la présentation de la logique applicative.

\item[Binding] Mécanisme JavaFX permettant de lier une propriété à une autre ou à une expression. Le modèle de données utilise \texttt{StringProperty}, \texttt{ListProperty} et \texttt{ObjectProperty} pour que les modifications dans l'interface se reflètent automatiquement dans le modèle et inversement.
\end{description}


\end{document}